\section{Host-memory tiers}
\label{sec:tiers}

A host tier moves the cliff by enlarging the effective capacity, and Eq.~\eqref{eq:restore} says it should cost little
to use. Whether it helps depends on whether it holds the working set, as Shao et~al.\ showed for
run time~\cite{shao2026efficientagent}; we measure how much energy it saves as a function of its size
(Table~\ref{tab:tiers}).

\begin{table}[t]
\centering\small
\caption{SGLang host tiers, Qwen3-32B, synthetic sessions, 5\,s pauses (two rounds). WT: write-through; WB: write-back. $N$: sessions; Hit: hit rate. GPU cache alone:
0.12\% hit and 1,058\,J per turn at 48 sessions, 0.09\% and 1,035\,J at 96. vs.\ GPU: change in J/turn against the GPU cache alone at the same load. Where the two rounds differ at the precision shown, both values are given.}
\label{tab:tiers}
\begin{tabular}{lrrrr}
\toprule
Tier & $N$ & Hit & J/turn & vs.\ GPU \\
\midrule
80\,GB WT & 48 & 11.4--12.6\% & 953--962 & $-10$ to $-9\%$ \\
80\,GB WB & 48 & 36.5\% & 768 & $-27\%$ \\
120\,GB WT & 48 & 49.9--50.3\% & 652--655 & $-38\%$ \\
160\,GB WT & 48 & 75.5\% & 450 & $-57\%$ \\
240\,GB WT & 48 & 75.5\% & 448--449 & $-58\%$ \\
160\,GB WT & 96 & 20.2\% & 873 & $-16\%$ \\
240\,GB WT & 96 & 37.2\% & 732 & $-29\%$ \\
\bottomrule
\end{tabular}
\end{table}

\paragraph{Size against the working set.}\enlargethispage{\baselineskip} At 48 sessions the working set is about 548k tokens. SGLang's write-through
tier mirrors the GPU, so its usable capacity is roughly the host size rather than the sum of host and
GPU~\cite{lu2026metastability}: 80\,GB holds about 305k tokens ($W/C_{\text{host}}\approx1.8$, past the cliff) and saved 9--10\%; 160\,GB
holds about 610k ($W/C_{\text{host}}\approx0.9$) and saved 57\% with 84\% more turns and a median time to first token of 9\,s
instead of 55\,s. A 240\,GB tier added nothing at 48 sessions, because the set already fit, and saved 29\% at 96
sessions, where it did not. vLLM's host tier, which we could size only to 70\,GiB in our container, changed energy by $+0.3\%$ to $-6\%$.

\paragraph{A less inclusive tier gets more from the same memory.} In write-back mode SGLang copies KV to the host only
when it is evicted from the GPU, so the two tiers overlap less. At 80\,GB, write-back raised the hit rate from 11.4--12.6\%
to 36.5\% and cut energy per turn by 27\% instead of 9--10\%; at 160\,GB it matched write-through (not shown). On recorded
trajectories at 192 and 288 sessions an 80\,GB write-back tier came within 1.3--4.5\% of 160\,GB write-through
(68.8 against 67.9\,J and 62.1 against 59.4\,J per turn) with turn counts within 4.5\%: a less inclusive tier with half the host memory gave nearly the same
result.

\paragraph{Real trajectories.} On recorded SWE-agent sessions the 160\,GB tier lowered energy per turn by 33\% with
33\% more turns at 96 sessions against SGLang's GPU cache alone (91.5 against 136.6--138.2\,J), and by 27\% and 31\%
with 22\% and 27\% more turns at 192 and 288 sessions against vLLM's GPU cache (we did not run SGLang without a
tier at these loads). At 48 sessions, with no memory pressure yet, it
saved 3\%. These gains are smaller than on synthetic sessions because the GPU cache alone did not
collapse (Section~\ref{sec:cliff}).

\paragraph{How well the one-parameter model does.} With $f=0.76$ fitted on the 160\,GB, 48-session cell,
Eq.~\eqref{eq:energy} predicts the 80\,GB tier at 962\,J (measured 958, two-round means), vLLM's 70\,GiB tier at 1,035 and 973\,J (measured 1,035
and 973), and every other overloaded (48- and 96-session) synthetic Qwen3-32B host-tier cell within 1.4\% (2.5--3.4\% at 16
sessions), except one: 160\,GB at 48 sessions with 30\,s pauses, where it
predicts 534\,J and we measured 614\,J ($-13\%$). That cell drew 563\,W on average against 578\,W with 5\,s pauses: idle time likely dilutes
the prefill share, which the model, built for a busy card, does not capture; Section~\ref{sec:production} shows the same limit on production traffic at light load. On Qwen3-8B it predicts 207 and 254\,J for the 160\,GB tier at
128 and 256 sessions, against 215 and 266\,J measured ($-4$ to $-5\%$). On recorded trajectories,
where the GPU cache never thrashed, with no cache taken as $E_{\text{thrash}}$, it predicts the 160\,GB tier 6--23\% below the
measured energy per turn, so the fit does not transfer to those workloads.

\paragraph{Host cost.} We could not read the host's energy counters. From the engines' CPU time (charged at 10\,W per
busy core) and the bytes moved (assuming 1\,nJ per byte of host memory traffic) we estimate, under these assumed coefficients, the host's extra energy, as a
share of the GPU energy saved, at 0.5--1.8\% for the 160\,GB tier at 48 and 96 sessions and 2.9--3.0\% for 80\,GB at 48;
where the GPU saving was small the same estimate is larger: 11--14\% for SGLang tiers at 16 sessions and for 80\,GB at
96, 128\% for 80\,GB at 16 sessions with 30\,s pauses, and 4.7--212\% for vLLM's 70\,GiB tier. Mao and Zhou bounded the host-side cost at below 1\% of the GPU saving on a
similar host~\cite{mao2026offloading}.
